



\begin{table*}[t]
    \centering
    \caption{
    Performance comparison of Qwen3 teacher and student across 4 benchmarks. The best results are in \textbf{bold}, and the second-best results are \underline{underlined}. We report average@32 over 5 runs.}
    \vspace{1 mm}
    \label{tab:main}

    \begingroup
    \setlength{\tabcolsep}{6.5 pt}
    \renewcommand{\arraystretch}{1.08}
    \normalsize

    \begin{adjustbox}{max width=0.82\textwidth}
    \begin{tabular}{llccccc}
        \toprule
        \multirow{2}{*}{\textbf{Params}} 
        & \multirow{2}{*}{\textbf{Method}} 
        & \multicolumn{2}{c}{\textbf{Math}} 
        & \textbf{Science} 
        & \textbf{Code} 
        & \multirow{2}{*}{\textbf{Average}} \\
        \cmidrule(lr){3-4}
        \cmidrule(lr){5-5}
        \cmidrule(lr){6-6}
        & & AIME 2024 & AIME 2025 & GPQA & LiveCodeBench & \\

        \midrule
        \rowcolor[HTML]{FEE090}
        \multicolumn{7}{c}{\textbf{Teacher Models from the Qwen3 Series}}\\
        \textbf{4B} 
        & GRPO 
        & 67.60\phantom{.}\scalebox{0.72}{$\pm$1.34} & 60.42\phantom{.}\scalebox{0.72}{$\pm$1.47} & 55.19\phantom{.}\scalebox{0.72}{$\pm$0.81} & 63.13\phantom{.}\scalebox{0.72}{$\pm$0.93} & 61.59 \\

        \midrule
        \rowcolor[HTML]{E0F3F8}
        \multicolumn{7}{c}{\textbf{Distilled Student Models from the Qwen3 Series}}\\
        \multirow{7}{*}{\textbf{0.6B}}
        & Vanilla 
        & 7.71\phantom{.}\scalebox{0.72}{$\pm$0.33} & 12.81\phantom{.}\scalebox{0.72}{$\pm$0.38} & 13.24\phantom{.}\scalebox{0.72}{$\pm$0.24} & 14.89\phantom{.}\scalebox{0.72}{$\pm$0.29} & 12.16 \\
        & SFT 
        & 5.67\phantom{.}\scalebox{0.72}{$\pm$0.26} & 5.42\phantom{.}\scalebox{0.72}{$\pm$0.22} & 15.20\phantom{.}\scalebox{0.72}{$\pm$0.30} & 9.61\phantom{.}\scalebox{0.72}{$\pm$0.27} & 8.97 \\
        & GRPO 
        & 4.06\phantom{.}\scalebox{0.72}{$\pm$0.37} & 4.90\phantom{.}\scalebox{0.72}{$\pm$0.31} & \underline{20.38}\phantom{.}\scalebox{0.72}{$\pm$0.47} & 15.95\phantom{.}\scalebox{0.72}{$\pm$0.51} & 11.32 \\
        & OPD 
        & \underline{16.82}\phantom{.}\scalebox{0.72}{$\pm$0.81} & \underline{22.95}\phantom{.}\scalebox{0.72}{$\pm$0.97} 
        & 17.76\phantom{.}\scalebox{0.72}{$\pm$0.41} 
        & \underline{22.65}\phantom{.}\scalebox{0.72}{$\pm$0.75} & \underline{20.04} \\
        & OPSD$_{\mathrm{gt}}$ 
        & 12.63\phantom{.}\scalebox{0.72}{$\pm$0.58} & 17.04\phantom{.}\scalebox{0.72}{$\pm$0.64} & 17.32\phantom{.}\scalebox{0.72}{$\pm$0.35} & 16.73\phantom{.}\scalebox{0.72}{$\pm$0.44} & 15.93 \\
        & OPSD$_{\mathrm{hint}}$ 
        & 9.77\phantom{.}\scalebox{0.72}{$\pm$0.43} & 14.12\phantom{.}\scalebox{0.72}{$\pm$0.48} & 15.98\phantom{.}\scalebox{0.72}{$\pm$0.30} & 12.65\phantom{.}\scalebox{0.72}{$\pm$0.37} & 13.13 \\
        & \textbf{\ourmodel} 
        & \textbf{20.84}\phantom{.}\scalebox{0.72}{$\pm$0.90} & \textbf{26.13}\phantom{.}\scalebox{0.72}{$\pm$1.07} & \textbf{22.19}\phantom{.}\scalebox{0.72}{$\pm$0.59} 
        & \textbf{27.72}\phantom{.}\scalebox{0.72}{$\pm$0.83} & \textbf{24.22} \\

        \midrule
        \multirow{7}{*}{\textbf{1.7B}}
        & Vanilla 
        & 9.90\phantom{.}\scalebox{0.72}{$\pm$0.38} & 8.96\phantom{.}\scalebox{0.72}{$\pm$0.33} & 26.80\phantom{.}\scalebox{0.72}{$\pm$0.36} & 22.73\phantom{.}\scalebox{0.72}{$\pm$0.41} & 17.10 \\
        & SFT 
        & 26.77\phantom{.}\scalebox{0.72}{$\pm$0.66} & 22.40\phantom{.}\scalebox{0.72}{$\pm$0.72} & 29.85\phantom{.}\scalebox{0.72}{$\pm$0.51} & 24.63\phantom{.}\scalebox{0.72}{$\pm$0.57} & 25.91 \\
        & GRPO 
        & 25.63\phantom{.}\scalebox{0.72}{$\pm$1.03} & 21.67\phantom{.}\scalebox{0.72}{$\pm$1.12} & 33.55\phantom{.}\scalebox{0.72}{$\pm$0.76} & 20.70\phantom{.}\scalebox{0.72}{$\pm$0.87} & 25.39 \\
        & OPD 
        & \underline{43.86}\phantom{.}\scalebox{0.72}{$\pm$1.23} & \underline{37.04}\phantom{.}\scalebox{0.72}{$\pm$1.31} 
        & 31.73\phantom{.}\scalebox{0.72}{$\pm$0.55} 
        & \underline{32.45}\phantom{.}\scalebox{0.72}{$\pm$0.93} 
        & \underline{36.27} \\
        & OPSD$_{\mathrm{gt}}$ 
        & 33.85\phantom{.}\scalebox{0.72}{$\pm$0.88} & 24.69\phantom{.}\scalebox{0.72}{$\pm$0.94} & \underline{35.02}\phantom{.}\scalebox{0.72}{$\pm$0.67} & 22.73\phantom{.}\scalebox{0.72}{$\pm$0.61} & 29.07 \\
        & OPSD$_{\mathrm{hint}}$ 
        & 34.42\phantom{.}\scalebox{0.72}{$\pm$0.76} & 21.43\phantom{.}\scalebox{0.72}{$\pm$0.83} & 33.46\phantom{.}\scalebox{0.72}{$\pm$0.59} & 23.12\phantom{.}\scalebox{0.72}{$\pm$0.52} & 28.11 \\
        & \textbf{\ourmodel} 
        & \textbf{50.83}\phantom{.}\scalebox{0.72}{$\pm$1.15} & \textbf{41.72}\phantom{.}\scalebox{0.72}{$\pm$1.24} & \textbf{38.72}\phantom{.}\scalebox{0.72}{$\pm$0.73} 
        & \textbf{40.63}\phantom{.}\scalebox{0.72}{$\pm$0.91} & \textbf{42.98} \\

        \bottomrule
    \end{tabular}%
    \end{adjustbox}

    \endgroup

    \vspace{-10pt}
\end{table*}



\subsection{Overall Performance}

We compare \ourmodel against six baselines on both 0.6B and 1.7B student models. The results are summarized in Table~\ref{tab:main}, from which we draw the following key observations. Note that we also provide training cost analysis in~\cref{app:cost}.

\begin{itemize}[leftmargin=*]

\item \textbf{Obs 1: \ourmodel consistently outperforms all baselines across every benchmark.}
On both model scales, \ourmodel achieves the highest scores on all four tasks, surpassing the strongest baseline OPD by \textbf{20.86\%} (0.6B) and \textbf{18.50\%} (1.7B) in relative average improvement. Notably, SFT and GRPO both fail to outperform the Vanilla baseline at the 0.6B scale, indicating that sparse outcome-level rewards and static demonstrations are insufficient for guiding small models in TIR. In contrast, our step-wise reweighting mechanism provides dense yet reliability-modulated supervision, effectively preventing misleading gradients from corrupted tool-call states while preserving informative teacher guidance in well-aligned regions.

\item \textbf{Obs 2: The improvements generalize across different model scales, demonstrating the scalability of \ourmodel.}
Our 1.7B student recovers \textbf{69.8\%} of the 4B teacher's performance, compared to only 58.9\% for OPD, showing that \ourmodel substantially improves distillation efficiency. Furthermore, even our 0.6B model surpasses several 1.7B baselines, suggesting that step-wise reweighting can partially compensate for the capacity gap between model scales. %

\end{itemize}




